% Copyright (c) 2026 Geovana Neves. Licensed under CC BY 4.0 (https://creativecommons.org/licenses/by/4.0/). See LICENSE-AND-AUTHORSHIP.md.
%
% 00-fts-overview/text/01-flightstream.tex
%
% Part 1: what FlightStream is, where it sits on the fidelity ladder, what it
% will not do, why the build is part of every answer, and the published
% record of the method. Manual pages are those of build 26121 [6].

\section{FlightStream}

\begin{frame}{Where a panel method sits}
\begin{columns}[T]
\begin{column}{0.46\textwidth}
\begin{itemize}\small
  \item Aerodynamic tools form a \textbf{fidelity ladder}: each level resolves
        more physics and costs more per case.
  \item A panel method sits in the middle: the \textbf{whole configuration} in
        three dimensions, in seconds to minutes per point on a workstation.
  \item That is the regime of \textbf{interference and installation}: where a
        propeller sits, what a fuselage does to a wing, a pitching moment slope.
\end{itemize}
\end{column}
\begin{column}{0.52\textwidth}
\slidefig[0.52]{g00-fidelity-ladder}
\end{column}
\end{columns}
\begin{takeaway}A panel method buys full-configuration aerodynamics at a cost
that still allows design iteration.\end{takeaway}
\end{frame}

\begin{frame}{What FlightStream is}
\begin{columns}[T]
\begin{column}{0.54\textwidth}
\begin{itemize}\small
  \item A \textbf{surface vorticity} panel solver: the potential-flow problem,
        with vorticity carried on the surface mesh [7, 8]. It tolerates
        unstructured, engineering-quality meshes and reads the common exchange
        formats [6, pp.~72-75].
  \item The core is \textbf{inviscid}. The rest is correction models on top of
        it: a boundary layer, with optional viscous coupling, for friction drag
        and decambering [6, pp.~204-211]; subsonic compressibility, chosen at
        initialisation [6, pp.~218-221], [12].
  \item Steady, unsteady and rotating solutions, driven from a graphical
        interface or from a \textbf{script} of commands.
\end{itemize}
\end{column}
\begin{column}{0.42\textwidth}
\begin{alertblock}{Out of scope by construction}
\begin{itemize}\small
  \item Massively separated flow, deep stall, buffet.
  \item Transonic shocks.
  \item Acoustics.
\end{itemize}
\end{alertblock}
\begin{alertblock}{And it does not warn you}
\small A case that leaves the validity envelope still converges and still
prints coefficients.
\end{alertblock}
\end{column}
\end{columns}
\end{frame}

\begin{frame}{The build number is part of the answer}
\begin{itemize}\small
  \item FlightStream is released as \textbf{builds}; this series names build
        \fsbuild{} in its examples. The scripting language changes between
        builds of one version: a command that runs on one build may be absent,
        renamed or spelled differently on another.
  \item The manual ships inside each build, and \textbf{editions are not page
        aligned}. A page number is quoted with its edition, or not at all.
  \item A coefficient without the build that produced it is \textbf{not
        reproducible}. Record the build next to every result.
\end{itemize}
\begin{block}{How pyflightstream holds this}
\small Every command it writes comes from a \textbf{command database per
build}, and every row of a study names its build in the column
\code{FS\_BUILD}. \pyfs{plan} checks the whole study against that build's
database before any seat is spent (Part 4, Guide 1).
\end{block}
\end{frame}

\begin{frame}{The tool's own view of the workflow}
\begin{columns}[T]
\begin{column}{0.50\textwidth}
\begin{itemize}\small
  \item The manual closes with three checklists: the \textbf{general}
        workflow [6, p.~411], the workflow from \textbf{CAD cross sections}
        (CCS) [6, p.~412], and the \textbf{rotor and turbine} workflow
        [6, p.~413].
  \item The general one has five steps: geometry preparation, half-model
        setup, surface detection, boundary conditions, simulation setup.
  \item Boundary conditions come \textbf{before} anything a user would call
        setup. Part 3 explains why: initialisation locks them.
\end{itemize}
\end{column}
\begin{column}{0.46\textwidth}
\begin{block}{A half model is expected}
\small The manual recommends the half model under a mirror plane, and free
edges along that plane are \textbf{expected}, not a defect [6, p.~411].
\end{block}
\begin{block}{A reference case ships with it}
\small A rotor in hover and cruise [6, p.~401]: the fastest check that an
installation behaves.
\end{block}
\end{column}
\end{columns}
\end{frame}

\begin{frame}{The method, and where it has been measured}
\begin{columns}[T]
\begin{column}{0.48\textwidth}
\begin{block}{What it solves}
\begin{itemize}\small
  \item The integrated-circulation formulation, in full [7].
  \item The surface vorticity panel method in the archival literature: the
        reference to cite when someone asks what FlightStream solves [8].
  \item The classical panel-method background: [9, 10].
\end{itemize}
\end{block}
\end{column}
\begin{column}{0.48\textwidth}
\begin{block}{Where it has been compared with experiment}
\begin{itemize}\small
  \item Stability and control of a vehicle against wind-tunnel data:
        derivatives agree acceptably, control effectiveness is overpredicted,
        propellers off [13].
  \item A high-lift wing with distributed propulsion, against experiment: the
        closest published case to a powered airframe [14].
\end{itemize}
\end{block}
\end{column}
\end{columns}
\vspace{1mm}
\begin{takeaway}Convergence against a theoretical anchor is
\emph{verification}; agreement with measured data is \emph{validation}. Say
which one a result shows.\end{takeaway}
\end{frame}
